\section{Nonlinear-risk indicators and ranking policy}
\label{sec:risk}

The ceilings of \cref{sec:freq} rank candidates by output capability. What
remains is the risk that a candidate, while formally inside its gates,
operates closer to its nonlinear regions than an alternative. Linear
correction does not help here: a linear pre-filter cannot reduce a
nonlinear residual for all programme material, because if it did, scaling
the input would force that residual to be linear.

This report does \emph{not} convert $\Xmax$, or any of the indicators
below, into an absolute distortion figure, and does not add unlike
indicators into a scalar total. They occur at different frequencies, with
different phases and different perceptual weights
\cite{klippel2006,allison1982,fielder2019}, so a sum would not correspond
to any measurable quantity. They are reported as a vector of separately
named, dimensionless indicators.

\subsection{Reported indicators}
\label{sec:risk-vector}

For each role evaluation the procedure reports

\begin{equation}
	\mathbf r=\big(\xi_x,\ \xi_x^{\mathrm{tr}},\ d_{\mathrm D},\
		\xi_{\mathrm{box}},\ \xi_P,\ \xi_{\mathrm{amp}},\
		f_{\mathcal B,\mathrm h}/f_L,\ V_{\mathrm b}^{\mathrm{tot}}\big) .
	\label{eq:riskvec}
\end{equation}

\textbf{Excursion utilization} $\xi_x=X_{\mathrm{dem}}/\Xmax$ and its
transient counterpart $\xi_x^{\mathrm{tr}}$ \cref{eq:xitr}. Expanding motor
force factor and suspension compliance about the rest position gives
even- and odd-order terms that grow with $\xi_x$, which is the motivation
for using utilization as a screening indicator; the coefficients are not
available from datasheets, so no distortion percentage is inferred.
$\xi_x=1$ is the hard clipping boundary (A7). The value $0.80$ is a
\emph{preferred} design margin---about $\SI{1.9}{dB}$ of excursion
headroom---used as a lexicographic preference and to size unit counts, not
as a rejection threshold and not as a distortion budget.

\textbf{Doppler sideband indicator.} For a role band $\mathcal B$, define
$X_{\mathcal B}=\max_{f\in\mathcal B}X_{\mathrm{dem}}(f)$ and let
$f_{\mathcal B,\mathrm h}$ be its upper edge. A component at
$f_{\mathcal B,\mathrm h}$ radiated from a cone moving with peak
$X_{\mathcal B}$ acquires the phase-modulation index
$m=2\pi f_{\mathcal B,\mathrm h}X_{\mathcal B}/c$. For small $m$, either
first-order sideband has amplitude ratio $J_1(m)/J_0(m)\simeq m/2$, so the
reported indicator is

\begin{equation}
	d_{\mathrm D}\equiv\frac{J_1(m)}{J_0(m)}
	\simeq\frac{\pi f_{\mathcal B,\mathrm h}X_{\mathcal B}}{c},
	\qquad
	X_{\mathcal B}=\max_{f\in\mathcal B}
		\frac{V_{\mathrm{dem}}(f)}{N\Sd} .
	\label{eq:doppler}
\end{equation}

The phase modulation is kinematic and exact within the rigid-piston model;
the sideband approximation is its small-index limit. At fixed output,
$d_{\mathrm D}$ falls as $1/\Sd$: among candidates with equal $\Vd$, the
larger cone has the lower indicator. It is \emph{not} a perceived
distortion value---the audibility of loudspeaker FM distortion has been
disputed since \cite{klipsch1969,allison1982}---and is used only as a
screen against a declared sideband-ratio budget. Splitting the bands
removes cross-band modulation \emph{on the same radiator}; it does not
remove self-modulation within a band, and finite crossover overlap leaves
some residual.

\textbf{Box-spring indicator.} Adiabatic compression of the enclosed air is
an asymmetric spring; weighting the leading term by the box's share of
total stiffness gives

\begin{equation}
	\xi_{\mathrm{box}} \approx
		\frac{\gamma+1}{4}\,
		\frac{V_{\mathrm{dem}}}{N\Vb}\Big(1-\frac{\Qts^2}{\Qtc^2}\Big),
	\label{eq:boxHD}
\end{equation}

a property of the enclosure rather than the motor, reduced by larger
$\Vb$ and by stuffing ($\gamma_{\mathrm{eff}}\to1$). It is reported
separately and never added to \cref{eq:doppler} or to $\xi_x$.

\textbf{Electrical and thermal utilization.} Let
\begin{equation}
	\begin{aligned}
	\xi_P &={P_{\mathrm{req}}}/{\Pmax},\\
	\xi_{\mathrm{amp}}&=\max\!\left\{
		\frac{\max(E_{\mathrm{req}},E_{\mathrm{burst}})}
		     {E_{\mathrm{amp}}},
		\frac{\max(I_{\mathrm{req}},I_{\mathrm{burst}})}
		     {I_{\mathrm{amp}}},\right.\\[-0.3ex]
	&\hspace{4.5em}\left.
		\frac{P_{\mathrm{req}}}{P_{\mathrm{amp,cont}}},
		\frac{P_{\mathrm{burst}}}{P_{\mathrm{amp,burst}}}
		\right\}.
	\end{aligned}
\end{equation}
These follow from \cref{eq:ampgates}. $\xi_P$ is a warning indicator only:
inferring coil temperature would require thermal resistances and time
constants that datasheets do not publish \cite{button1992}, and $\Pmax$
itself carries standard-dependent uncertainty (A5).

\textbf{Inductive screen.} $f_{\mathcal B,\mathrm h}/f_L$ indicates how far the role's upper
band edge intrudes into the region where $\Le$ and its position dependence
matter. Since $\Le$ is frequency- and level-dependent, $f_L$ is a
screening approximation rather than a model-validity boundary (A2). The
declared policy nevertheless rejects a reported $f_L\le f_2$; a missing
$\Le$ is recorded as unresolved rather than imputed, and is demoted in the
upper-role order.

\textbf{Enclosure size.} The total box volume $V_{\mathrm b}^{\mathrm{tot}}$
of a role is reported and used as a late tie-break, because it is a real
constraint of the application rather than a risk.

\subsection{Ranking policy}
\label{sec:risk-policy}

The policy is deliberately non-compensatory and is declared in advance:

\begin{enumerate}
\item \textbf{Hard gates.} $\xi_x\le1$, $\xi_x^{\mathrm{tr}}\le1$,
	\cref{eq:ampgates}, $\Qtc\le\Qtc^\ast$ achievable within the enclosure
	limit, $d_{\mathrm D}$ within budget for the upper role, and
	$f_L>f_{\mathcal B,\mathrm h}$ when $\Le$ is reported. Failure at any
	band frequency rejects the candidate, with the
	binding gate reported.
\item \textbf{Pareto front.} Within a role, surviving candidates are sorted
	into non-dominated fronts on the role's core objectives:
	$\max(\xi_x,\xi_x^{\mathrm{tr}})$, $\xi_{\mathrm{amp}}$, and box volume
	for the low-bass role, with $d_{\mathrm D}$ added for the upper-bass
	role. Restricting the front to the core objectives avoids the degenerate
	outcome in which almost every candidate is non-dominated; the full
	vector \cref{eq:riskvec} remains reported.
\item \textbf{Lexicographic order within a front.} The low-bass role
	minimizes excursion utilization, then Doppler sideband ratio, then
	electrical demand, then box-spring indicator, then total volume. The upper-bass
	role first prefers candidates at or below the $0.80$ excursion margin,
	then maximizes $\eta_0\Pmax$, then prefers a higher $f_L$, then
	minimizes excursion and electrical demand.
\item \textbf{Pairing.} A system candidate is a feasible low-bass manifold
	plus a feasible upper-bass channel. Roles are compared by their own keys
	in sequence, then by physical driver count and total enclosure volume.
	Risk indicators of different roles are never summed.
\end{enumerate}

The lexicographic order is an engineering policy, not a derived optimum.
Its influence on the outcome is therefore tested explicitly by re-running
the ranking under alternative reasonable orders
(\cref{sec:procedure-robustness}).
